\subsection*{Chapter 13: Pretraining on Unlabeled Data}

\paragraph{Exercise 13.1 (Initial loss).}
The loss at one position is $-z_y+\log\sum_u e^{z_u}$. Since $z_y$ has mean zero,
$\E[\mathcal L]=\E[\log\sum_u e^{z_u}]$. For large $V$ the sum concentrates around its mean,
$V\,\E[e^{z}]=V e^{\sigma^2/2}$ (log-normal mean), so $\E[\mathcal L]\approx\ln V+\sigma^2/2$.
For GPT-2 small, the final LayerNorm makes each hidden coordinate roughly unit variance and the
head weights have standard deviation $0.02$, so each logit has variance
$\approx768\cdot0.02^2=0.31$. Prediction: $\ln 50{,}257+0.15\approx10.98$ nats.

\paragraph{Exercise 13.2 (Perplexity across tokenizers).}
Model A: total NLL $=3.10\times10^6$ nats over $4.3\times10^6$ bytes $=0.721$ nats/byte $=1.04$
bits/byte. Model B is byte-level, so tokens $=$ bytes and $0.95$ nats/byte $=1.37$ bits/byte.
Model A is better, even though its per-token loss is three times larger.

\paragraph{Exercise 13.3 (Training vs.\ validation).}
Expect the validation minimum within the first few epochs; dropout delays it slightly and
keeps the validation loss lower afterwards, because it limits how precisely the model can store
the training windows. Samples near the validation minimum are novel but story-like; late samples
reproduce training passages verbatim, especially with greedy decoding.

\paragraph{Exercise 13.4 (Temperature limits).}
Let $z_1>z_j$ for $j\neq1$. Then $p_\tau(j)/p_\tau(1)=e^{(z_j-z_1)/\tau}\to0$ as $\tau\to0^+$,
so $p_\tau\to e_1$. As $\tau\to\infty$, $z_j/\tau\to0$ for all $j$ and $p_\tau(j)\to1/V$. With
$k$ tied maxima the limit is uniform over the tied set (the ratios between tied tokens are always
1), which differs from \code{argmax}, which deterministically returns the first index.

\paragraph{Exercise 13.5 (Do top-$k$ and temperature commute?).}
Yes for top-$k$: dividing by $\tau>0$ preserves the order of the logits, so the same $k$ tokens
survive, and the softmax over the survivors is the same either way. No for top-$p$: the kept set
depends on the cumulative probabilities, which temperature changes. A high temperature flattens
the distribution, so more tokens are needed to reach mass $p$. Applying temperature first (the
common convention) lets temperature widen the nucleus.

\paragraph{Exercise 13.6 (Resume correctness).}
Required: model weights; optimizer state (moments, step counts); scheduler/step counter; CPU and
CUDA RNG states (dropout); Python \code{random} and NumPy RNGs if used anywhere; the sampler's
generator state and the position within the epoch (how many batches have been consumed); worker
seeding (\code{worker\_init\_fn}/base seed) so augmentation inside workers replays; and the
\code{GradScaler} state (current scale and growth tracker). Listing~13.5 omits the data-loader
position and sampler RNG, the Python/NumPy RNGs and the \code{GradScaler}; autocast itself holds
no state.

\paragraph{Exercise 13.7 (Checkpoint interval).}
Job MTBF $M=50{,}000\,\text{h}/2048=24.4\,\text{h}\approx87{,}900\,\text{s}$.
Synchronous ($\delta=90$\,s): $T=\sqrt{2\cdot90\cdot87{,}900}\approx3{,}980$\,s ($\approx66$ min).
Expected overhead $\approx\delta/T+T/(2M)\approx2.3\%+2.3\%=4.5\%$.
Asynchronous ($\delta=6$\,s): $T\approx1{,}030$\,s ($\approx17$ min), overhead
$\approx0.6\%+0.6\%=1.2\%$ (restart time excluded in both).

\paragraph{Exercise 13.8 (Accumulation is not always exact).}
(1) Padding: if each micro-batch's loss is the mean over its \emph{non-padding} tokens and the
micro-batches contain different numbers of real tokens, averaging the $k$ means weights tokens
unequally; the exact large-batch loss divides the summed NLL by the total token count. Fix: sum
token losses and divide by the global token count. (2) Layers with batch statistics (BatchNorm)
see smaller batches. Also: the last, partial micro-batch of an epoch, and dropout masks that
differ because the RNG is consumed in a different pattern.

\paragraph{Exercise 13.9 (Activation memory, GPT-2 xl).}
Model states: $16\times1.558\text{B}=24.9$\,GB; budget $0.9\cdot80-24.9=47.1$\,GB.
(a) No FlashAttention: per sample $48\cdot1024\cdot1600\cdot(34+5\cdot25\cdot1024/1600)=8.97$\,GB
$\Rightarrow b=5$. (b) FlashAttention: $48\cdot1024\cdot1600\cdot34=2.67$\,GB $\Rightarrow b=17$.
(c) Full checkpointing: block inputs $2\cdot1024\cdot1600\cdot48=0.157$\,GB plus one layer's full
activations during recomputation ($0.187$\,GB) $\approx0.344$\,GB per sample $\Rightarrow b\approx136$
(in practice temporary buffers and the logits tensor limit this much earlier).

\paragraph{Exercise 13.10 (Measure MFU).}
Typical ordering: fp32 lowest (no tensor cores; TF32 if enabled helps); bf16 autocast a large
jump; SDPA a further gain (no $T\times T$ materialization, fewer memory-bound kernels); compile
another 10--30\% from fusing LayerNorm/GELU/residual chains and reducing launches; activation
checkpointing reduces MFU by roughly a quarter (recompute is not counted in MFU) but allows larger
micro-batches, which can partly win it back.

\paragraph{Exercise 13.11 (Chinchilla planning).}
$C=1000\cdot86{,}400\cdot0.4\cdot989\times10^{12}=3.4\times10^{22}$ FLOPs, so
$N_{\text{opt}}=\sqrt{C/120}\approx16.9$B and $D_{\text{opt}}\approx338$B tokens.
Inference over a year: $1.2\times10^{13}$ tokens $\times2N$. Option A ($N=16.9$B):
$4.1\times10^{23}$ inference FLOPs, total $\approx4.4\times10^{23}$. Option B ($N=8.4$B,
$D=1.35$T): training $6\cdot8.4\text{B}\cdot1.35\text{T}=6.8\times10^{22}$ (2$\times$ the budget),
inference $2.0\times10^{23}$, total $\approx2.7\times10^{23}$. Option B is about 40\% cheaper in total;
with the Chinchilla loss fit its loss is only slightly higher than option A's, which is the
argument for over-training models that will be served heavily.

\paragraph{Exercise 13.12 (Chunked loss).}
Implement an \code{autograd.Function} whose forward takes the final hidden states $\mat H$
\code{(N, d)}, the head weight $\mat W$ \code{(V, d)} and targets. Loop over chunks of rows: compute
$\mat Z_c=\mat H_c\mat W^\top$ in fp32, its log-sum-exp, the chunk loss, and immediately the
gradient $\partial\mathcal L/\partial\mat Z_c=(\softmax(\mat Z_c)-\text{onehot})/N$; accumulate
$\partial\mathcal L/\partial\mat H_c=\partial\mathcal L/\partial\mat Z_c\,\mat W$ and
$\partial\mathcal L/\partial\mat W\mathrel{+}=(\partial\mathcal L/\partial\mat Z_c)^\top\mat H_c$.
Save only these gradients for backward (which just scales them by the incoming gradient). Peak
memory drops from $N\cdot V$ to $\text{chunk}\cdot V$ floats; verify with
\code{torch.autograd.gradcheck} on small sizes in float64.
